% !TEX root = ../matan.tex

\section{Теорема Банаха о сжимающем отображении}

\subsection*{Сжимающие отображения}

\begin{Definition}{Сжимающее отображение}{}
	Пусть $(X,\rho)$ --- метрическое пространство. Отображение $F\colon X\to X$
	называется \textbf{сжимающим}, если существует константа $0\le c<1$ такая, что
	\[
	\rho\bigl(F(x),F(y)\bigr)\le c\,\rho(x,y)\qquad\text{для всех } x,y\in X.
	\]
	Число $c$ называется коэффициентом сжатия. Всякое сжимающее отображение
	непрерывно (даже равномерно непрерывно).
\end{Definition}

\subsection*{Теорема о неподвижной точке}

\begin{Theorem}{Банаха о неподвижной точке}{}
	Если $(X,\rho)$ --- полное метрическое пространство, а $F\colon X\to X$ ---
	сжимающее отображение, то $F$ имеет единственную неподвижную точку, то есть
	единственную точку $x_0\in X$ с $F(x_0)=x_0$.
\end{Theorem}

\begin{proof}
	\textbf{Существование.} Зафиксируем произвольную точку и построим итерационную
	последовательность
	\[
	x_1\in X,\quad x_2=F(x_1),\quad x_3=F(x_2),\quad\dots,\quad x_{n+1}=F(x_n).
	\]
	Оценим расстояние между соседними членами. Применяя свойство сжатия $k-1$ раз,
	\[
	\rho(x_k,x_{k+1})=\rho\bigl(F(x_{k-1}),F(x_k)\bigr)\le c\,\rho(x_{k-1},x_k)
	\le\dots\le c^{k-1}\rho(x_1,x_2).
	\]
	Пусть $m>n$. По неравенству треугольника и сумме геометрической прогрессии,
	\[
	\rho(x_n,x_m)\le\sum_{k=n}^{m-1}\rho(x_k,x_{k+1})
	\le\rho(x_1,x_2)\sum_{k=n}^{m-1}c^{k-1}
	\le\rho(x_1,x_2)\,\frac{c^{n-1}}{1-c}.
	\]
	Так как $0\le c<1$, правая часть стремится к нулю при $n\to\infty$ (равномерно по
	$m>n$). Следовательно, для любого $\eps>0$ найдётся $N$ такое, что при $m>n\ge N$
	выполнено $\rho(x_n,x_m)<\eps$, то есть $\{x_n\}$ --- последовательность Коши. В
	силу полноты $X$ существует предел $x_0=\lim_{n\to\infty}x_n$.

	Перейдём к пределу в равенстве $x_{n+1}=F(x_n)$. Левая часть стремится к $x_0$, а
	правая --- к $F(x_0)$ ввиду непрерывности $F$. Поэтому $x_0=F(x_0)$, то есть $x_0$
	--- неподвижная точка.

	\textbf{Единственность.} Пусть $x_0$ и $x_0'$ --- две неподвижные точки. Тогда
	\[
	\rho(x_0,x_0')=\rho\bigl(F(x_0),F(x_0')\bigr)\le c\,\rho(x_0,x_0').
	\]
	Отсюда $(1-c)\,\rho(x_0,x_0')\le 0$. Так как $1-c>0$, получаем
	$\rho(x_0,x_0')=0$, то есть $x_0=x_0'$.
\end{proof}
